\documentclass[11pt,oneside]{article}

\pdfminorversion=7
\usepackage{QuizStyle}

\hypersetup{
    pdftitle={20261022 Quiz for MATH301 Modern Algebra I},
    pdfauthor={Jungwoo Kim},
    pdfsubject={Quiz for MATH301 Modern Algebra I},
    pdfcreator={LaTeX}
}

\begin{document}

\quiztitle{20261022 Quiz}

\begin{exercise}{4.5.13}
Prove that a group of order $56$ has a normal Sylow $p$-subgroup for some prime $p$ dividing its order.
\end{exercise}

\begin{solution}
Since $56=2^3\cdot7$, Sylow's Theorem gives
\[
    n_7\mid8,\qquad n_7\equiv1\pmod7,
\]
so $n_7=1$ or $8$. If $n_7=1$, the Sylow $7$-subgroup is normal.

If $n_7=8$, the eight distinct subgroups of order $7$ intersect pairwise only in the identity and account for $8(7-1)=48$ elements of order $7$. Thus exactly eight elements of $G$ do not have order $7$. Every Sylow $2$-subgroup has order $8$ and contains no element of order $7$, so it consists of precisely these eight elements. Hence the Sylow $2$-subgroup is unique and therefore normal.
\end{solution}

\begin{exercise}{4.5.16}
Let $|G|=pqr$, where $p,q,r$ are primes with $p<q<r$. Prove that $G$ has a normal Sylow subgroup for $p$, $q$, or $r$.
\end{exercise}

\begin{solution}
Sylow's Theorem gives $n_r\mid pq$ and $n_r\equiv1\pmod r$. Since $p,q<r$, this forces $n_r=1$ or $pq$. In the first case the Sylow $r$-subgroup is normal.

Suppose that $n_r=pq$. These distinct subgroups of order $r$ account for $pq(r-1)$ nonidentity elements, leaving only $pq-1$ nonidentity elements of other orders. On the other hand, $n_q\mid pr$ and $n_q\equiv1\pmod q$. If $n_q>1$, then $n_q\ge r$, since $p<q$; the Sylow $q$-subgroups would consequently contribute at least $r(q-1)$ distinct elements of order $q$. But
\[
    r(q-1)\ge rp>pq-1,
\]
a contradiction. Hence $n_q=1$, and the Sylow $q$-subgroup is normal.
\end{solution}

\begin{exercise}{4.5.17}
Prove that if $|G|=105$, then $G$ has a normal Sylow $5$-subgroup and a normal Sylow $7$-subgroup.
\end{exercise}

\begin{solution}
Since $105=3\cdot5\cdot7$, Sylow's Theorem gives
\[
    n_5\in\{1,21\},\qquad n_7\in\{1,15\}.
\]
If both were greater than $1$, the distinct Sylow $5$- and Sylow $7$-subgroups would contain $21(5-1)+15(7-1)=174$ distinct nonidentity elements, which is impossible. Thus at least one is normal.

\Needspace{15\baselineskip}
Suppose first that a Sylow $7$-subgroup $P_7$ is normal. The quotient $G/P_7$ has order $15$ and hence has a unique Sylow $5$-subgroup. Its preimage $K$ in $G$ is a normal subgroup of order $35$. In $K$, the number of Sylow $5$-subgroups divides $7$ and is congruent to $1$ modulo $5$, so it equals $1$. Conjugation by any element of $G$ preserves $K$ and therefore its unique Sylow $5$-subgroup. Thus a Sylow $5$-subgroup of $G$ is normal.

Conversely, if a Sylow $5$-subgroup $P_5$ is normal, then $G/P_5$ has order $21$ and a unique Sylow $7$-subgroup. Its preimage is a normal subgroup $K$ of order $35$; in $K$, the Sylow $7$-subgroup is unique, since its number divides $5$ and is $1$ modulo $7$. The same conjugation argument shows that it is normal in $G$.
\end{solution}

\begin{supplementaryexercise}
Let $G$ be a group of order $30$. Determine whether $G$ is simple. Justify your answer.
\end{supplementaryexercise}

\begin{solution}
Sylow's Theorem gives
\[
    n_5\in\{1,6\},\qquad n_3\in\{1,10\}.
\]
If $G$ were simple, neither Sylow subgroup could be normal, so $n_5=6$ and $n_3=10$. The distinct Sylow $5$-subgroups would then contain $6(5-1)=24$ elements of order $5$, while the distinct Sylow $3$-subgroups would contain $10(3-1)=20$ elements of order $3$. Including the identity gives at least $1+24+20=45$ elements, contradicting $|G|=30$. Hence $G$ is not simple.
\end{solution}

\begin{exercise}{5.1.18}
In each of (a) to (e), give an example of a group with the specified properties:
\begin{exercisesubparts}
    \item an infinite group in which every element has order $1$ or $2$;
    \item an infinite group in which every element has finite order but for each positive integer $n$ there is an element of order $n$;
    \item a group with an element of infinite order and an element of order $2$;
    \item a group $G$ such that every finite group is isomorphic to some subgroup of $G$;
    \item a nontrivial group $G$ such that $G\cong G\times G$.
\end{exercisesubparts}
\end{exercise}

\begin{solution}
\begin{exercisesubparts}
    \item Take $G=\bigoplus_{n\ge1}C_2$. It is infinite, and every element has finite support and squares to the identity; thus every nonidentity element has order $2$.

    \item Take $G=\bigoplus_{n\ge1}C_n$. Every element has only finitely many nonidentity coordinates, so its order is finite. The generator in the $n$th factor has order $n$ for each $n\ge1$.

    \item Take $G=\Z\times C_2$. The element $(1,1_{C_2})$ has infinite order, while $(0,t)$ has order $2$ when $t$ is the nonidentity element of $C_2$.

    \item Take $G=\bigcup_{n\ge1}S_n$, where $S_n\le S_{n+1}$ fixes $n+1$. Every finite group $H$ embeds into $S_{|H|}$ by its action on itself by left multiplication (Cayley's Theorem), and hence embeds into $G$.

    \item Take $G=\bigoplus_{n\ge1}C_2$ again. Splitting the coordinates into odd and even indices gives $G\cong G\times G$. The group is nontrivial.
\end{exercisesubparts}
\end{solution}

\Needspace{24\baselineskip}
\begin{exercise}{5.2.10}
Let $n$ and $k$ be positive integers and let $A$ be the free abelian group of rank $n$ (written additively). Prove that $A/kA$ is isomorphic to the direct product of $n$ copies of $\Z/k\Z$, where
\[
    kA=\{ka\mid a\in A\}.
\]
\end{exercise}

\begin{hint}
See Exercise 5.1.14.
\end{hint}

\begin{solution}
Choose a basis $a_1,\ldots,a_n$ of $A$. Each $a\in A$ is uniquely of the form $a=m_1a_1+\cdots+m_na_n$ with $m_i\in\Z$. Define
\[
\begin{aligned}
    \varphi\colon A&\longrightarrow(\Z/k\Z)^n,\\
    m_1a_1+\cdots+m_na_n&\longmapsto([m_1]_k,\ldots,[m_n]_k).
\end{aligned}
\]
This is a surjective homomorphism. Moreover, $a\in\ker\varphi$ if and only if $k\mid m_i$ for every $i$, which is equivalent to $a\in kA$. Thus $\ker\varphi=kA$, and the First Isomorphism Theorem gives
\[
    A/kA\cong(\Z/k\Z)^n.
\]
\end{solution}

\end{document}
